\subsection{Objetivos de diseño}

El diseño de la arquitectura se orientó por tres objetivos principales. El primero fue preservar la exactitud numérica de la simulación, de modo
que la representación comprimida no introdujera diferencias respecto a una ejecución de referencia no comprimida. El segundo consistió en reducir
la memoria ocupada por el vector de estado sin exigir su reconstrucción completa ante cada operación. El tercero fue mantener una sobrecarga
temporal razonable, de forma que el ahorro espacial no dependiera de una penalización prohibitiva en tiempo de ejecución.

Estos objetivos implican que la arquitectura debía cumplir dos condiciones simultáneamente. Por un lado, toda transformación aplicada al
almacenamiento del estado debía ser estrictamente reversible. Por otro, el acceso al vector debía adaptarse al patrón local con que las compuertas
actúan sobre las amplitudes, evitando tratamientos globales cuando la actualización solo involucra una fracción del arreglo.

Desde esta perspectiva, el diseño no busca sustituir el modelo de simulación de estado completo, sino modificar su representación física en
memoria. El vector de amplitudes se mantiene como referencia lógica del simulador, pero su materialización pasa a depender de una capa intermedia
capaz de alternar entre forma comprimida y forma densa según las necesidades de acceso de cada operación. Esta separación entre representación
lógica y representación física es la base que articula el resto de la propuesta.
